\documentclass[12pt]{article}
\usepackage[utf8]{inputenc}
\usepackage[T1]{fontenc}
\usepackage{hyperref}
\title{Ukhwejo Language Resources for K-12 Curriculum Development}
\author{Compiled by Claude Sonnet 5 (claude-sonnet-5) for Andr\'as Kornai}
\date{}
\begin{document}
\maketitle
\section{Introduction}
Ukhwejo (also Benkonjo; ISO 639-3 ukh) is a Bantu language of the Makaa-Njem group (Guthrie A.802), spoken by about 2,000 people (1996 figure) in scattered villages in south-western Central African Republic. Sango is the wider language of communication. Verbix lists the language as threatened and unwritten. Joshua Project gives 3,300 for the Benkonjo people and notes that Scripture translation is still needed.
\section{Primary Educational Resources}
No educational materials were found in the search results.
\section{Linguistic and Reference Materials}
\begin{itemize}
\item Wikipedia, Ukhwejo language: \url{https://en.wikipedia.org/wiki/Ukhwejo_language}
\item Abstract by Christina Thornell (WOCAL) on the cluster's location and linguistic identity: \url{https://www.uni-koeln.de/phil-fak/afrikanistik/wocal/schedule/abstracts/2-5-3%20Christina%20Thornell.pdf}
\item Verbix Endangered Languages entry: \url{https://docs.verbix.com/EndangeredLanguages/Ukhwejo}
\end{itemize}
\section{Religious and Community Resources}
\begin{itemize}
\item Joshua Project profile of Benkonjo in Central African Republic: \url{https://joshuaproject.net/people_groups/10794}
\end{itemize}
\section{Digital and Audio Resources}
\begin{itemize}
\item Endangered Languages Project entry: \url{https://www.endangeredlanguages.com/lang/565?hl=en}
\end{itemize}
\section{Curriculum Development Context}
The language is unwritten and thinly documented, with no teaching materials found. Any curriculum would first need an orthography and community-led documentation, with Sango as the bridge language.
\end{document}